% !TeX root = cheat_sheet.tex

\documentclass[../final.tex]{subfiles}

\begin{document}

\setcounter{chapter}{-1}
\chapter{Вспомогательная литература}


% ============================================================
% СИСТЕМЫ СЧИСЛЕНИЯ
% ============================================================

\section{Системы счисления: основные обозначения}

Число в системе с основанием $q$:

\[
(a_n a_{n-1}\ldots a_1a_0.a_{-1}a_{-2}\ldots a_{-m})_q.
\]

\[
(a_n\ldots a_0.a_{-1}\ldots a_{-m})_q
=
\sum_{k=-m}^{n} a_kq^k,
\qquad
0\leq a_k<q.
\]

\medskip

\begin{center}
\renewcommand{\arraystretch}{1.25}
\begin{tabular}{cccc}
\textcolor{rc3}{\bfseries Основание}
&
\textcolor{rc3}{\bfseries Система}
&
\textcolor{rc3}{\bfseries Цифры}
&
\textcolor{rc3}{\bfseries Пример}
\\[2mm]

$2$
&
двоичная
&
$0,1$
&
$101101_2$
\\

$8$
&
восьмеричная
&
$0,\ldots,7$
&
$572_8$
\\

$10$
&
десятичная
&
$0,\ldots,9$
&
$935_{10}$
\\

$16$
&
шестнадцатеричная
&
$0,\ldots,9,A,\ldots,F$
&
$D6_{16}$
\end{tabular}
\end{center}

\[
A=10,\qquad
B=11,\qquad
C=12,\qquad
D=13,\qquad
E=14,\qquad
F=15.
\]


% ============================================================

\section{Десятичная система счисления}

\subsection{Перевод \texorpdfstring{$q \to 10$}{q → 10}}

Каждая цифра умножается на соответствующую степень основания:

\[
(a_n\ldots a_0.a_{-1}\ldots a_{-m})_q
=
a_nq^n+\ldots+a_1q+a_0
+a_{-1}q^{-1}
+\ldots
+a_{-m}q^{-m}.
\]

\subsubsection{\texorpdfstring{$2 \to 10$}{2 → 10}}

\[
101101.01_2
\]

\[
=
1\cdot2^5
+0\cdot2^4
+1\cdot2^3
+1\cdot2^2
+0\cdot2^1
+1\cdot2^0
+0\cdot2^{-1}
+1\cdot2^{-2}
\]

\[
=
32+8+4+1+\frac14
=
45.25_{10}.
\]

\subsubsection{\texorpdfstring{$8 \to 10$}{8 → 10}}

\[
572_8
=
5\cdot8^2+7\cdot8+2
\]

\[
=
320+56+2
=
378_{10}.
\]

\subsubsection{\texorpdfstring{$16 \to 10$}{16 → 10}}

\[
3A7_{16}
=
3\cdot16^2+10\cdot16+7
\]

\[
=
768+160+7
=
935_{10}.
\]


% ============================================================

\subsection{Перевод \texorpdfstring{$10 \to q$}{10 → q}}

\subsubsection{Целая часть}

Последовательное деление на $q$:

\[
\begin{aligned}
N   &=qN_1+r_0,\\
N_1 &=qN_2+r_1,\\
N_2 &=qN_3+r_2,\\
    &\ \vdots
\end{aligned}
\qquad
0\leq r_i<q.
\]

Результат:

\[
N_{10}
=
(r_k r_{k-1}\ldots r_1r_0)_q.
\]

\textcolor{rc3}{\bfseries Остатки читаются снизу вверх.}

\medskip

Пример:

\[
45_{10}\longrightarrow ?_2
\]

\[
\begin{array}{rcl}
45 &=& 2\cdot22+1,\\
22 &=& 2\cdot11+0,\\
11 &=& 2\cdot5+1,\\
5  &=& 2\cdot2+1,\\
2  &=& 2\cdot1+0,\\
1  &=& 2\cdot0+1.
\end{array}
\]

\[
45_{10}=101101_2.
\]


\subsubsection{Дробная часть}

Последовательное умножение на $q$:

\[
qx_i=a_{-(i+1)}+x_{i+1},
\qquad
0\leq x_{i+1}<1.
\]

\textcolor{rc3}{\bfseries Целые части читаются сверху вниз.}

\medskip

Пример:

\[
0.625_{10}\longrightarrow ?_2
\]

\[
\begin{array}{rcl}
0.625\cdot2 &=& 1.25,\\
0.25 \cdot2 &=& 0.50,\\
0.50 \cdot2 &=& 1.00.
\end{array}
\]

\[
0.625_{10}=0.101_2.
\]


\subsubsection{Целая и дробная части}

Переводятся отдельно:

\[
45.625_{10}
=
45_{10}+0.625_{10}.
\]

\[
45_{10}=101101_2,
\qquad
0.625_{10}=0.101_2.
\]

\[
45.625_{10}=101101.101_2.
\]


% ============================================================

\section{Перевод \texorpdfstring{$q_1 \to q_2$}{q1 → q2}}

\subsection{Перевод \texorpdfstring{$2 \to 8$}{2 → 8}}

Разбить двоичное число на группы по $3$ бита,
начиная от точки.

\[
\begin{array}{cccccccc}
000&001&010&011&100&101&110&111\\[1mm]
0&1&2&3&4&5&6&7
\end{array}
\]

\subsubsection{Целое число}

\[
11010110_2
\]

Дополняем слева нулём:

\[
011\;010\;110
\]

\[
011\longrightarrow3,
\qquad
010\longrightarrow2,
\qquad
110\longrightarrow6.
\]

\[
11010110_2=326_8.
\]


\subsubsection{Число с дробной частью}

Группировка идёт от точки в обе стороны:

\[
10110.1011_2
\]

\[
010\;110
\;.\;
101\;100
\]

\[
010\longrightarrow2,
\qquad
110\longrightarrow6,
\]

\[
101\longrightarrow5,
\qquad
100\longrightarrow4.
\]

\[
10110.1011_2=26.54_8.
\]


% ============================================================

\subsection{Перевод \texorpdfstring{$8 \to 2$}{8 → 2}}

Каждую восьмеричную цифру заменить соответствующей тройкой битов:

\[
\begin{array}{cccccccc}
0&1&2&3&4&5&6&7\\[1mm]
000&001&010&011&100&101&110&111
\end{array}
\]

Пример:

\[
572_8
\]

\[
5\longrightarrow101,
\qquad
7\longrightarrow111,
\qquad
2\longrightarrow010.
\]

\[
572_8
=
101\;111\;010_2.
\]

\[
572_8=101111010_2.
\]


% ============================================================

\subsection{Перевод \texorpdfstring{$2 \to 16$}{2 → 16}}

Разбить двоичное число на группы по $4$ бита,
начиная от точки.

\medskip

\begin{center}
\renewcommand{\arraystretch}{1.2}
\begin{tabular}{cccccccc}
$0000$ & $0001$ & $0010$ & $0011$ &
$0100$ & $0101$ & $0110$ & $0111$ \\

$0$ & $1$ & $2$ & $3$ &
$4$ & $5$ & $6$ & $7$
\end{tabular}

\vspace{2mm}

\begin{tabular}{cccccccc}
$1000$ & $1001$ & $1010$ & $1011$ &
$1100$ & $1101$ & $1110$ & $1111$ \\

$8$ & $9$ & $A$ & $B$ &
$C$ & $D$ & $E$ & $F$
\end{tabular}
\end{center}


\subsubsection{Целое число}

\[
101111100011_2
\]

\[
1011\;1110\;0011
\]

\[
1011\longrightarrow B,
\qquad
1110\longrightarrow E,
\qquad
0011\longrightarrow3.
\]

\[
101111100011_2=BE3_{16}.
\]


\subsubsection{Число с дробной частью}

\[
101101.1101_2
\]

\[
0010\;1101
\;.\;
1101
\]

\[
0010\longrightarrow2,
\qquad
1101\longrightarrow D.
\]

\[
101101.1101_2=2D.D_{16}.
\]


% ============================================================

\subsection{Перевод \texorpdfstring{$16 \to 2$}{16 → 2}}

Каждую шестнадцатеричную цифру заменить соответствующей
четвёркой битов.

\medskip

\begin{center}
\renewcommand{\arraystretch}{1.2}
\begin{tabular}{cccccccc}
$0$ & $1$ & $2$ & $3$ &
$4$ & $5$ & $6$ & $7$ \\

$0000$ & $0001$ & $0010$ & $0011$ &
$0100$ & $0101$ & $0110$ & $0111$
\end{tabular}

\vspace{2mm}

\begin{tabular}{cccccccc}
$8$ & $9$ & $A$ & $B$ &
$C$ & $D$ & $E$ & $F$ \\

$1000$ & $1001$ & $1010$ & $1011$ &
$1100$ & $1101$ & $1110$ & $1111$
\end{tabular}
\end{center}

Пример:

\[
D6_{16}
\]

\[
D\longrightarrow1101,
\qquad
6\longrightarrow0110.
\]

\[
D6_{16}
=
1101\;0110_2.
\]

\[
D6_{16}=11010110_2.
\]


% ============================================================

\subsection{Перевод \texorpdfstring{$8 \leftrightarrow 16$}{8 ↔ 16}}

Перевод выполняется через двоичную систему.


\subsubsection{\texorpdfstring{$8 \to 16$}{8 → 16}}

\[
326_8
\]

Сначала:

\[
3\longrightarrow011,
\qquad
2\longrightarrow010,
\qquad
6\longrightarrow110.
\]

\[
326_8
=
011\;010\;110_2.
\]

Перегруппируем по $4$ бита:

\[
011010110_2
=
0000\;1101\;0110_2.
\]

\[
1101\longrightarrow D,
\qquad
0110\longrightarrow6.
\]

\[
326_8=D6_{16}.
\]


\subsubsection{\texorpdfstring{$16 \to 8$}{16 → 8}}

\[
D6_{16}
\]

Сначала:

\[
D\longrightarrow1101,
\qquad
6\longrightarrow0110.
\]

\[
D6_{16}
=
1101\;0110_2.
\]

Перегруппируем по $3$ бита:

\[
11010110_2
=
011\;010\;110_2.
\]

\[
011\longrightarrow3,
\qquad
010\longrightarrow2,
\qquad
110\longrightarrow6.
\]

\[
D6_{16}=326_8.
\]

\newpage

% ============================================================

\section{Дополнительные материалы}

\subsection{Быстрая таблица перевода}

\begin{center}

\renewcommand{\arraystretch}{1.25}

\begin{tabular}{cccc}
\textcolor{rc3}{\bfseries $10$}
&
\textcolor{rc3}{\bfseries $2$}
&
\textcolor{rc3}{\bfseries $8$}
&
\textcolor{rc3}{\bfseries $16$}
\\[2mm]

0  & 0000 & 0  & 0 \\
1  & 0001 & 1  & 1 \\
2  & 0010 & 2  & 2 \\
3  & 0011 & 3  & 3 \\
4  & 0100 & 4  & 4 \\
5  & 0101 & 5  & 5 \\
6  & 0110 & 6  & 6 \\
7  & 0111 & 7  & 7 \\
8  & 1000 & 10 & 8 \\
9  & 1001 & 11 & 9 \\
10 & 1010 & 12 & A \\
11 & 1011 & 13 & B \\
12 & 1100 & 14 & C \\
13 & 1101 & 15 & D \\
14 & 1110 & 16 & E \\
15 & 1111 & 17 & F

\end{tabular}

\end{center}


% ============================================================

\subsection{Краткий алгоритм выбора перевода}

\[
\begin{array}{ccl}
q\rightarrow10
&:&
\displaystyle
\sum_k a_kq^k,
\\[3mm]

10\rightarrow q
&:&
\text{деление / умножение на }q,
\\[3mm]

2\rightarrow8
&:&
\text{группы по }3\text{ бита},
\\[2mm]

8\rightarrow2
&:&
\text{каждая цифра }\rightarrow3\text{ бита},
\\[2mm]

2\rightarrow16
&:&
\text{группы по }4\text{ бита},
\\[2mm]

16\rightarrow2
&:&
\text{каждая цифра }\rightarrow4\text{ бита},
\\[2mm]

8\leftrightarrow16
&:&
\text{через двоичную систему}.
\end{array}
\]


% ============================================================

\subsection{Полезные степени двойки}

\[
\begin{array}{c|ccccccccc}
n
&0&1&2&3&4&5&6&7&8
\\
2^n
&1&2&4&8&16&32&64&128&256
\end{array}
\]

\[
\begin{array}{c|cccccccc}
n
&9&10&11&12&13&14&15&16
\\
2^n
&512&1024&2048&4096&8192&16384&32768&65536
\end{array}
\]

\[
2^{10}=1024\approx10^3,
\qquad
2^{20}\approx10^6,
\qquad
2^{30}\approx10^9.
\]


\end{document}
